\section{正規分布と推測統計学}
\label{s08_CorrelationCausation}

\subsection{授業内容}
\subsubsection{科目の中でのこのコマの位置づけ}
\begin{tabular}[t]{ccccccc}
    記述統計[3] & 確率[2] & モデル[1] & 推測・推定[1] & 意思決定[0] & 実践[1]
\end{tabular}

\subsubsection{概要}
前回に引き続き，連続的な分布の例として正規分布を導入する。正規分布は様々な影響が混合している現象の結果として得られる変数の形であり，心理学を始め多くの学問で基本的な分布として扱われる。二項分布のような離散変数の確率分布とは違い，連続的な確率分布になるため確率密度で表現されることに注意が必要である。誤差分布が正規分布に従うと考える場合，その平均は0であり分散だけが興味の対象になる。どのパラメタが形状にどのような変化を与えるのか，視覚的イメージとして理解しておくことも重要である。正規分布に従うデータについては，理論分布を使って相対的な比較が可能であるが，その際は素点ではなく標準化されたスコアを使い，標準正規分布に照らして考える。
\subsubsection{コマ主題細目}
\begin{description}
\item[正規分布の形]心理統計学において最もよく使われる確率分布が正規分布である。二項分布の極限としても表現されるため，誤差を表す分布とも考えられる。正規分布は連続変数であり，平均値と分散がそれぞれ位置母数，形状母数として与えられる。それぞれの母数が変化すると正規分布の形(Sahpe)がどのように変わるか，また形状の変化にかかわらず一定範囲の中に含まれる相対度数が定められていることを理解する。
\begin{flushright}
    $\to$ \citet{野島201903}のP.68--70，\citet{Yamada2004}のP.80--85など，おそらく心理統計のテキストであれば大抵，正規分布には言及している。\citet{長沼201611}のP38--39も参照。
\end{flushright}

\item[正規分布の関数]正規分布の形状は，ガウスが発見した確率密度関数で描かれる。正規分布の確率密度関数は，複雑な数式をしているが，その特徴は左右対象で単峰であることにある。確率密度関数がどうしてこのような形になっているのか，関数を読み解きながら確率分布の形を理解する。
\begin{flushright}
    $\to$ \citet{長沼201611}のP.62--69.
\end{flushright}

\item[連続変数の確率分布関数]離散変数の確率分布とは違い，連続変数における確率密度関数は，密度が高さであり確率そのものではないことに注意が必要である。変数が連続であるとはどういうことかを正しく理解し，その上で確率密度がどのように定義されるか，また確率密度関数では確率が面積として表されることを理解する。数学的な詳細には立ち入らないが，数学的側面について直観的に理解しておく必要がある
\begin{flushright}
    $\to$ \citet{Maeda2017}のP.84--88.
\end{flushright}

\item[標準正規分布]様々な現象が正規分布に従うとしても，平均値や標準偏差が異なるのであれば相対的な比較はしにくい。しかし素点を標準化しておくと，標準正規分布に照らし合わせることで推測をすることができる。数字の表や統計環境Rをもちいた例を示しつつ，標準正規分布の活用の仕方について理解する。
\begin{flushright}
    $\to$ \citet{川端201408}のP.36--44，\citet{Yamada2004}のP.86--89。
\end{flushright}

\end{description}
\subsubsection{キーワード}
\begin{itemize}
\item 正規分布，$\mu$と$\sigma$
\item 標準正規分布
\end{itemize}
\subsection{授業情報}
\paragraph{コマの展開方法}
講義
\paragraph{標準シラバスにおける位置づけ}
科目番号5 心理学統計法；1.心理学で用いられる統計手法；E 連続分布：正規分布を中心として 
\subsubsection{予習・復習課題}
\paragraph{予習}
二項分布の計算をしっかり復習した後で，「パスカルの三角形」や「ゴルトンボード」について事前に調べておくと良い。後者については特に動画でみておくことで正規分布についての直観的な理解を得ることができる。
\paragraph{復習}
正規分布の数式は大変複雑な形をしているように見えるが，その本質は「左右対称」かつ「単峰」であることである。この点については\citet{長沼201611}が言葉による解説を行なっており，数式に苦手意識を持っていてもわかりやすいので一読するとよい。連続変数における確率は確率密度として表現されることについて，\citet{Maeda2017}のP.84--88を読みながら復習すると良い。最後に身近なデータをつかって，変数の標準化と標準正規分布における相対的な位置の推測の仕方について，確認しておこう。計算に際してはExcel/Numbers/Calcsなどの表計算ソフトが持っている確率密度関数や，統計環境Rのそれを使うとよい。